% Số: facts.md (Test v3 và Test-3: 32,9 -> 1,2%, 32,4 -> 3,8%; 40 ca lọt v1 theo nguyên nhân, mục "Bổ sung sau phản
% biện 3"; vi mô broker, mục "Bổ sung cho bài v3"). Giai đoạn đầu (luật chuỗi trên hai tập niêm phong) được kể như một
% phần của chính nghiên cứu này, không gọi là phiên bản trước.
\section{Giới thiệu}\label{sec:intro}

Trong mô hình ReAct~\cite{yao2023react}, tác tử LLM xen kẽ suy luận với lời gọi công cụ, và nội dung công cụ trả về đi
thẳng vào ngữ cảnh của mô hình. Trong tấn công chèn lệnh gián tiếp~\cite{greshake2023indirect}, kẻ tấn công cài chỉ thị
vào nội dung mà tác tử sẽ đọc để khiến nó gửi dữ liệu tới một đích lạ, làm lộ giá trị mật hoặc sửa nội dung cần gửi.
Phòng thủ mức prompt phụ thuộc vào việc mô hình có tuân theo hay không, còn các bộ phân loại chèn lệnh phần lớn được
huấn luyện trên tiếng Anh. Một hướng khác đặt việc cưỡng chế ở phía host, giữa mô hình và công cụ, theo nguyên lý
reference monitor~\cite{anderson1972}.

Cách đơn giản nhất để cưỡng chế tại Tool Broker là so khớp chuỗi. Broker chỉ thực thi một lời gọi gửi ra ngoài khi đích
của nó là một chuỗi con của yêu cầu người dùng $u$, và chỉ cho một giá trị lấy từ nguồn gắn nhãn mật rời tác tử khi $u$
chứa nguyên văn giá trị đó. Ở giai đoạn đầu của nghiên cứu này, các luật như vậy kết hợp với lời nhắc kiểu sandwich đã
giảm tỉ lệ tấn công thành công trên hai tập kiểm thử tiếng Việt niêm phong, với ba mô hình nền mở, từ 32,9\% xuống
1,2\% và từ 32,4\% xuống 3,8\%. Phân tích cả 40 ca còn lọt cho thấy ba chỗ hở (Bảng~\ref{tab:mechanisms}). Có 25 ca làm
lộ một giá trị mật nằm trong nguồn gắn nhãn công khai, vì nhãn được đặt cho cả nguồn. Có 5 ca gửi tới một đích có xuất
hiện trong $u$ nhưng không được cấp quyền, chẳng hạn đích nằm trong một câu cấm. Mười ca còn lại sửa tham số của một lời
gọi tới đích hợp lệ.

Các ca này cho thấy một chuỗi xuất hiện trong $u$ chưa chắc đã được cấp quyền, và nhãn của nguồn không cho biết độ nhạy
cảm của từng giá trị. Bài báo đề xuất \ours, một cơ chế cưỡng chế tại broker gồm ba thành phần tất định nhắm vào ba chỗ
hở này. \POne{} trích từ $u$ các quyền có cấu trúc, xét phạm vi phủ định và loại trừ, và chỉ cấp quyền cho đích khớp
nguyên token. \PTwo{} gán độ nhạy cảm cho từng giá trị trong đầu ra công cụ theo kiểu của giá trị, độc lập với nhãn
nguồn. \PThree{} buộc lời gọi được cấp quyền giữ các giá trị người dùng đã nêu. Chính sách chỉ được suy ra từ $u$,
không dựa vào phán đoán của một LLM, và cho kết quả \emph{cần xác nhận} khi quyền không rõ ràng. Bài báo có ba đóng góp.
\begin{enumerate}
\item \textbf{Phương pháp.} \ours{} xử lý ba chỗ hở đã đo được của luật so khớp chuỗi mà không gọi thêm LLM, với thời
  gian kiểm dưới một mili giây mỗi lời gọi (Mục~\ref{sec:method}).
\item \textbf{Mô hình hình thức.} Ba mệnh đề về cấp quyền, bảo mật và toàn vẹn được chứng minh dưới năm giả định, tách
  khỏi các thuộc tính chỉ đo được bằng thực nghiệm (Mục~\ref{sec:formal}).
\item \textbf{Đánh giá.} \ours{} được đóng băng rồi chạy một lần trên một tập kiểm thử mới gồm 48 họ, niêm phong trước
  khi có mã, với giả thuyết đăng ký trước, ablation, hai vòng tấn công thích nghi và AgentDojo bằng hai ngôn ngữ
  (Mục~\ref{sec:protocol}, \ref{sec:results}). Bằng chứng mức ca cho thấy mỗi thành phần chặn đúng loại tấn công mà
  nó nhắm tới, còn khẳng định thống kê so với luật so khớp chuỗi cần một tập kiểm thử lớn hơn
  (Mục~\ref{sec:discussion}).
\end{enumerate}
Dữ liệu, mã và các manifest niêm phong được công bố công khai.
